\chapter{كيف تتعرف على قطة}
% full version: chapters/12_recognition.tex

\pencilfig{12}{\pencilart{12}}{قطة كبيرة وقطة صغيرة: صورتان مختلفتان تمامًا على مستشعرات العين، والجواب واحد.}

قد تكون القطة في الصورة يسارًا أو يمينًا، قريبة أو بعيدة، في الشمس أو في الظل. في كل مرة يسقط على مستشعرات العين نمط مختلف تمامًا، ومع ذلك تتعرف عليها في جزء من الثانية. وهنا تكمن صعوبة التعرف: يجب أن نتعرف على الشيء نفسه في آلاف الهيئات المختلفة.

\section*{خط التجميع}

يحل الدماغ المسألة بخط تجميع من نحو عشر مراحل، وتقطعه الإشارة في عُشر ونصف من الثانية. في كل مرحلة يجري أمران. أولًا ”و“: يعمل العصبون إذا كانت القطع المطلوبة في الترتيب المطلوب، كأن يلتقي قطعتان مستقيمتان في زاوية. ثم ”أو على كل المواضع“: يعمل العصبون التالي إذا وُجدت الزاوية في أي مكان قريب، لا يهم أين بالضبط. العملية الأولى تجعل التعرف دقيق التمييز، والثانية تجعله متسامحًا مع الإزاحة. مرحلة بعد مرحلة تتجمع القطع في أجزاء، والأجزاء في أشياء، وتُغتفر الإزاحات وتغيرات الحجم.

\pencilfig{112}{%
\foreach \i/\y/\cx/\cy in {1/1.2/0.15/0.3,2/0.3/0.3/0.1,3/-0.6/0.05/0.05}{
  \draw[penlite] (0,\y-0.3) rectangle (0.6,\y+0.3);
  \draw[pcGray, line width=0.8pt] (\cx,\y-0.3+\cy+0.35) -- (\cx,\y-0.3+\cy) -- (\cx+0.3,\y-0.3+\cy);
  \draw[pen=pcBlue, wash=pcBlue] (1.9,\y) circle (0.2);
  \draw[arr] (0.65,\y) -- (1.65,\y);
  \draw[arr=pcBlue] (2.12,\y) -- (3.62,{0.3+0.12*(2-\i)});}
\draw[pen=pcBlue, wash=pcBlue] (3.9,0.3) circle (0.25);
\draw[arr] (4.2,0.3) -- (5.75,0.3); \node[lbl, anchor=south] at (4.97,0.35) {مراحل أخرى};
% a small cat
\draw[penlite] (6.3,0.3) circle (0.32);
\draw[penlite] (6.03,0.5) -- (6.07,0.82) -- (6.23,0.6); \draw[penlite] (6.57,0.5) -- (6.53,0.82) -- (6.37,0.6);
\node[lbl, anchor=north] at (0.3,-1.0) {الصورة}; \node[lbl, anchor=north, align=center] at (1.9,-1.0) {و: زاوية\\في هذا الموضع};
\node[lbl, anchor=north, align=center] at (3.9,-1.0) {أو: زاوية\\في أي موضع}; \node[lbl, anchor=north] at (6.3,-1.0) {قطة};
}{مرحلة واحدة من خط التجميع: ”و“ تجمع الزاوية في موضع واحد، و”أو“ تغتفر موضعها بالضبط.}

هذا المخطط بالذات نسخته الشبكات العصبية الالتفافية التي تتعرف على الصور في هاتفك. وهي تردّ الجميل: الشبكات المدرَّبة على التعرف على الأشياء تتنبأ جيدًا إلى حد ما بكيفية استجابة العصبونات في المراحل العليا من الإبصار.

\section*{التعلم من الفيديو}

لكن من علّم الدماغ؟ لم يعرض عليك أحد مليون صورة موسومة ”قطة“. جواب معقول: الفيديو. العالم يتدفق بلا انقطاع، وفي جزء من الثانية لا يكاد الشيء يتغير، بل يتغير موضعه وإضاءته. إذا قويت الوصلات بين ما يُرى الآن وما كان قبل لحظة، ربطت الشبكة بنفسها الهيئات المختلفة للشيء الواحد: البطيء هو الشيء، والسريع هو الظروف. وثمة تجربة استُبدلت فيها الأشياء خلسة في أثناء حركات العين، فأخذت العصبونات تعد أشياء مختلفة شيئًا واحدًا. أما بوصفه قاعدة عامة فلا يزال هذا فرضية.

\section*{الخدع البصرية بصمات للآليات}

الخدع البصرية ليست أعطالًا، بل آثار لطريقة بناء الآلة. إذا نظرت العين طويلًا إلى شلال ثم إلى حجارة ساكنة، ”زحفت“ الحجارة إلى الأعلى: زوج القناتين ”أسفل“ و”أعلى“ اختل توازنه بسبب تعب إحداهما. والفستان الشهير الذي يراه بعضهم أبيض وذهبيًا وآخرون أزرق وأسود يقسم الناس بحسب الإضاءة التي يفترضها دماغهم بصمت. و”الأفاعي الدوّارة“ المرسومة تبدو متحركة لأن المناطق الساطعة تُعالج أسرع قليلًا من الباهتة، ويُقرأ فرق الزمن حركةً. أما المكعب الذي يبرز تارة ويغوص تارة فهو مجموعتان من العصبونات تتنافسان؛ ووفق أحد النماذج، الضجيج هو ما يحسم تبدّل النتيجة في الغالب.

والطريف أن شبكة عصبية مدرَّبة على التنبؤ بالإطار التالي في الفيديو ”ترى“ هي أيضًا الأفاعي تدور. إذن بعض الخدع البصرية ثمن لا مفر منه للقدرة على التنبؤ.

\section*{الدماغ في مواجهة الشبكة العصبية}

التشابه ليس تامًا. تكفي الدماغ أمثلة قليلة، ويتعلم بلا وسوم تقريبًا، وينفق على كل شيء عشرين واطًا. أما الشبكة العصبية فيمكن خداعها بتزوير في الصورة لا يلحظه الإنسان. لكن تزويرًا كهذا يربك البشر أيضًا قليلًا إذا عُرض عليهم لحظة. أين تمر الحدود الحقيقية بين هاتين الآلتين؟ هذا ما لا يزال قيد البحث.

\fullversion{12}
